
The integral over $R_3$ goes to zero by strict optimality of $\thetatrue$
and the proper behavior of the prior.

\begin{lem}\lemlabel{region_r3}
%
Under \assuref{core_bclt_assu}, when $\regevent$ occurs with
$\epsilon_U < \prior(\thetatrue) / 2$,
%
\begin{align*}
%
\norm{I_3(\phi)}_2 =
O(\exp(-N \epsilon_L)) \expect{\prior(\theta)}{\norm{ \phi(\theta, \xvec_s) }_2 }.
%
\end{align*}
%
\begin{proof}
%
We have
%
\begin{align*}
\norm{I_3(\phi)}_2
={}& \norm{ \int_{R_3} \phi(\theta, \xvec_s)
       \exp\left(
            N \likhat(\theta) - N \likhat(\thetahat)
        \right)
       d \theta }_2 \\
\le&
\sup_{\theta \in R_3}
\exp\Bigg(
    N \Bigg(
         \frac{1}{N} \sumn \ell(\x_n \vert \theta) -
         \frac{1}{N} \sumn \ell(\x_n \vert \thetahat)
     \Bigg)
 \Bigg) \times
 \int_{R_3} \norm{\phi(\theta, \xvec_s)}_2
 \frac{\prior(\theta)}{\prior(\thetahat)} d \theta \\
={}&
    \exp(-N \epsilon_L)
    \frac{1}{\prior(\thetahat)}
    \expect{\prior(\theta)}{
        \norm{\phi(\theta, \xvec_s)}_2 \mathbb{I}(\theta \in R_3)
    }
    \quad\textrm{(\lemref{regular_event}, \eventref{strict_opt})}\\
\le&
    \exp(-N \epsilon_L)
    \frac{1}{\prior(\thetatrue) - \epsilon_U}
    \expect{\prior(\theta)}{\norm{\phi(\theta, \xvec_s)}_2 \mathbb{I}(\theta \in R_3)}
\quad\textrm{(\lemref{regular_event}, \eventref{prior_close})}\\
\le&
\exp(-N \epsilon_L) 2 \prior(\thetatrue)^{-1}
\expect{\prior(\theta)}{\norm{\phi(\theta, \xvec_s)}_2 }.
\end{align*}
%
In the last step, we take $\epsilon_U < \prior(\thetatrue) / 2$ in $\regevent$
(specifically, \lemref{regular_event} \eventref{prior_close}).
%
\end{proof}
%
\end{lem}
%
We note that the form of the bound in \lemref{region_r3} is designed to
take advantage of the prior mean assumption \defref{bclt_okay} \itemref{prior_op1},
or its more general version \defref{bclt_okay_index} \itemref{prior_op1_index}.
